\section{Geometry and trace formulas over K}\label{sec:geometry}

Fix compatible primitive cube roots in
\(K=\Q(\zeta_3)\) and \(E=\Q_2(\zeta_3)\), and put
\(\mathcal O=\Z_2[\zeta_3]\),
\(S=\operatorname{Spec}\mathcal O_K[1/6]\), and
\(U=\PP^1_S\setminus\{0,1,\infty\}\).
Let \(\pi:\mathcal E\to U\) be the smooth projective Legendre elliptic
scheme, with its zero section and affine equation
\[
 \mathcal E_\lambda:\quad y^2=x(x-1)(x-\lambda),
 \qquad \mathcal V=R^1\pi_*E.
\]
The cubic torsor \(f:T\to U\), \(z^3=\lambda\), has deck
transformation \(a(z)=\zeta_3z\). We use pullback as the deck action.
Let \(\mathcal L_\rho\) be the \(\zeta_3\)-eigensheaf of \(a^*\)
in \(f_*E\). This definition fixes the arithmetic normalization;
in particular no constant Galois twist is inserted. For \(r\in\{2,4\}\)
set
\[
 \mathcal G_r=\Sym^r\mathcal V\otimes\mathcal L_\rho,
 \qquad
 H_r=H^1_!(U_{\bar K},\mathcal G_r)
 :=\operatorname{im}\bigl(H_c^1(U_{\bar K},\mathcal G_r)
                  \longrightarrow H^1(U_{\bar K},\mathcal G_r)\bigr).
\]
Thus \(W_2=H_2\) and \(V_4=H_4\). The subscript records the symmetric
power, and is distinct from the order three of the Kummer character.

All Frobenius operators below are geometric. Write \(\tau\) for the
character of \(E(-1)\), so \(\tau(\Frob_{\mathfrak p})=N\mathfrak p\).
For a prime \(\mathfrak p\nmid6\), put \(q=N\mathfrak p\) and let
\(\rho_q:\F_q^\times\to E^\times\) be the trace character of the
specified eigensheaf. For even \(k\) define
\[
 a_q(\lambda)=q+1-\#\mathcal E_\lambda(\F_q),\qquad
 M_{3,k}(q)=\sum_{\lambda\in\F_q\setminus\{0,1\}}
                  \rho_q(\lambda)a_q(\lambda)^k.
\]
The singular parameters are excluded throughout these moment sums.

\begin{theorem}\label{thm:common-geometry}
For \(r=2,4\), the representation \(H_r\) is continuous of dimension
\(r\) over \(E\), and is unramified outside the primes of \(K\) above
\(2,3\). There is an exact sequence of \(G_K=\Gal(\bar K/K)\)-representations
\begin{equation}\label{eq:common-full-sequence}
 0\longrightarrow E\longrightarrow H_c^1(U_{\bar K},\mathcal G_r)
 \longrightarrow H^1(U_{\bar K},\mathcal G_r)
 \longrightarrow E(-(r+1))\longrightarrow0.
\end{equation}
Both middle spaces have dimension \(r+1\), and their common parabolic
image is \(H_r\). At every \(\mathfrak p\nmid6\),
\begin{align}
 \Tr(\Frob_{\mathfrak p}\mid W_2)
   &=-M_{3,2}(q)-q-1,\label{eq:trace-W2}\\
 \Tr(\Frob_{\mathfrak p}\mid V_4)
   &=-M_{3,4}(q)+3qM_{3,2}(q)+q^2-1.
   \label{eq:trace-V4}
\end{align}
These assertions include the degree-two prime ideals of \(K\).
\end{theorem}

\subsection{Local monodromy and the arithmetic boundary}

At \(\lambda=0,1\), the discriminant
\(16\lambda^2(1-\lambda)^2\) has valuation two, while
\(c_4=16(\lambda^2-\lambda+1)\) is a unit. Thus the fibres are
multiplicative, with nontrivial unipotent inertia on \(\mathcal V\).
At infinity, the substitutions \(t=1/\lambda\), \(x=\lambda X\),
\(y=\lambda^2Y\) give
\[
 Y^2=tX(X-1)(X-t).
\]
This is the ramified quadratic twist of the multiplicative family at
\(t=0\). The geometric inertia types on \(\mathcal V\) at
\(0,1,\infty\) are consequently \(J_2(1),J_2(1),J_2(-1)\);
see also \cite[Chapter~15]{Katz}.

For clarity, the symmetric-power calculation is made over the
characteristic-zero field \(E\). In a local basis choose
\(N(e)=0\), \(N(h)=ce\), \(c\ne0\), for the unipotent logarithm.
The induced logarithm satisfies
\[
 N_r(e^{r-j}h^j)=jc\,e^{r-j+1}h^{j-1},\qquad
 N_r^r(h^r)=r!c^re^r\ne0.
\]
It therefore has one block of length \(r+1\). Since \(r\) is even,
the scalar \(-1\) at infinity disappears. No assertion that \(r!\)
is an integral unit is needed. With compatible tame generators the
Kummer scalars at zero and infinity are a primitive cube root \(\eta\)
and \(\eta^{-1}\), respectively. Hence
\begin{center}
\begin{tabular}{c|c|c|c}
 point & \(\mathcal V\) & \(\mathcal G_r\) &
          \(\dim_E\mathcal G_r^{I}\)\\ \hline
 \(0\) & \(J_2(1)\) & \(\eta J_{r+1}(1)\) & \(0\)\\
 \(1\) & \(J_2(1)\) & \(J_{r+1}(1)\) & \(1\)\\
 \(\infty\) & \(J_2(-1)\) & \(\eta^{-1}J_{r+1}(1)\) & \(0\)
\end{tabular}
\end{center}
The same calculation holds in every good geometric fibre. The integral
tameness argument below verifies that no Swan terms occur there.

\begin{lemma}\label{lem:arithmetic-boundary}
The inertia invariant and coinvariant lines at \(\lambda=1\) are
arithmetically
\[
 \mathcal G_r^{I_1}\simeq E,
 \qquad (\mathcal G_r)_{I_1}\simeq E(-r).
\]
In particular geometric Frobenius acts on the invariant line by \(1\),
both over \(K\) and at every good residue field.
\end{lemma}

\begin{proof}
Put \(v=x-1\) in the singular fibre at \(\lambda=1\). Its equation
\(y^2=(1+v)v^2\) has the two rational tangents \(y=\pm v\), so its
multiplicative reduction is split. The Tate torsion sequence, followed by
inverse limit and duality, gives the cohomological sequence
\[
 0\longrightarrow E\longrightarrow\mathcal V
      \longrightarrow E(-1)\longrightarrow0
\]
on the local punctured trait; see \cite[pp.~18--19]{Tate}.
Its nontrivial unipotent inertia fixes the first line. In its \(r\)-th
symmetric power, the invariant line is the \(r\)-th power of this
trivial line, whereas the coinvariant quotient is the \(r\)-th power
of the quotient \(E(-1)\). Finally the torsor fibre over \(1\) is
\(\mu_3\). It is rational over \(K\) and over each residue field at
\(\mathfrak p\nmid6\), so its chosen character line is arithmetically
trivial. Tensoring by that line gives both assertions.
\end{proof}

\subsection{Cohomology and the two boundary terms}

In this subsection cohomology is first taken over \(\bar K\); the same
arguments apply over good algebraically closed residue fields.
The absence of \(I_0\)-invariants implies
\(H^0(U,\mathcal G_r)=0\). It also implies
\(H^0(U,\mathcal G_r^\vee(1))=0\), since a constant Tate twist is
trivial on geometric inertia. A lisse section on a connected nonproper
smooth curve cannot have proper support unless it is zero. Poincare
duality \cite[Theorem~24.1]{Milne} therefore gives
\begin{equation}\label{eq:common-vanishing}
 H_c^0(U,\mathcal G_r)=H_c^2(U,\mathcal G_r)=0.
\end{equation}
Comparison with the complex thrice-punctured sphere, which retracts
onto a graph with one vertex and two loops, gives ordinary Euler
characteristic \(-(r+1)\) for a rank-\((r+1)\) local system.
Duality gives the same compact-support Euler characteristic. The
specialization below transports the dimensions to the good fibres.
In particular
\(\dim_E H_c^1(U,\mathcal G_r)=r+1\).

Let \(j:U_{\bar K}\hookrightarrow\PP^1_{\bar K}\), and let
\(i_1\) be the inclusion of the point \(1\). Lemma~\ref{lem:arithmetic-boundary}
and the local invariant table give an exact sequence of ordinary etale
sheaves
\[
 0\longrightarrow j_!\mathcal G_r\longrightarrow j_*\mathcal G_r
     \longrightarrow i_{1*}E\longrightarrow0.
\]
Its cohomology sequence and the vanishing of \(H^0(U,\mathcal G_r)\)
give
\begin{equation}\label{eq:common-left-boundary}
 0\longrightarrow E\longrightarrow H_c^1(U,\mathcal G_r)
   \longrightarrow H^1(\PP^1,j_*\mathcal G_r)\longrightarrow0.
\end{equation}
The edge map for Leray injects
\(H^1(\PP^1,j_*\mathcal G_r)\) into \(H^1(U,\mathcal G_r)\).
The composite of this injection with
\eqref{eq:common-left-boundary} is the natural compact-support map.
It follows that
\[
 H^1(\PP^1,j_*\mathcal G_r)=H_r,
 \qquad \dim_EH_r=r.
\]
These identifications and exact sequences are Galois equivariant.

To determine the right boundary, the same sheaf sequence and
\eqref{eq:common-vanishing} imply
\(H^2(\PP^1,j_*\mathcal G_r)=0\). Leray now gives a surjection from
\(H^1(U,\mathcal G_r)\) onto
\(\bigoplus_{s=0,1,\infty}H^1(I_s,\mathcal G_r)\), with kernel
\(H_r\). At \(0,\infty\), the operator \(T_s-1\) is invertible,
so the corresponding local groups vanish. At \(1\), tame inertia
cohomology identifies the local group with
\[
 H^1(I_1,\mathcal G_r)
   \simeq (\mathcal G_r)_{I_1}(-1)
   \simeq E(-(r+1)).
\]
The additional twist comes from the dual of the tame factor
\(\Z_2(1)\); geometric Frobenius acts on this factor of the local
cohomology by \(q\). Combining these observations proves
\eqref{eq:common-full-sequence}, including its final surjectivity, and
\(\dim_EH^1(U,\mathcal G_r)=r+1\).

\subsection{Integral tame specialization}

We give the specialization argument explicitly because \(U\) is not
proper. The eigensheaf has an integral lattice
\(\mathcal L_{\rho,\mathcal O}\): the order-three projector is
integral since \(3\) is invertible in \(\mathcal O\). Put
\[
 \mathcal G_{r,\mathcal O}
  =\bigl(\Sym^r(R^1\pi_*\Z_2)\otimes_{\Z_2}\mathcal O\bigr)
       \otimes_{\mathcal O}\mathcal L_{\rho,\mathcal O},
 \qquad \mathcal G_{r,m}=\mathcal G_{r,\mathcal O}/2^m.
\]
The finite monodromy of \(\mathcal G_{r,m}\) factors through the image
of a subgroup of
\(\operatorname{GL}_2(\Z/2^m)\times\mu_3\), whose order is
\[
 18\cdot2^{4(m-1)}.
\]
Thus for every residue characteristic \(p\ge5\), wild pro-\(p\)
inertia at a boundary point acts trivially. In characteristic zero there
is no wild inertia. This verifies relative tameness in the fibrewise
discrete-valuation sense of \cite[Expose~XIII, Definition~2.1.1]{SGA1}.
The sheaves are already lisse on \(U\).

Apply \cite[Expose~XIII, Corollary~2.9]{SGA1} to
\(\PP^1_S\), its relative normal-crossing boundary consisting of the
three disjoint sections, and the additive coefficient sheaf
\(\mathcal G_{r,m}\). The compactification is smooth, proper, and of
finite presentation. The coefficient sheaf is finite, locally constant,
constructible, relatively tame, and \(2\)-primary, with \(2\)
invertible on \(S\). The theorem supplies local constancy and
base change in degrees at most one. The same argument applies to the
integral dual coefficient systems and their Tate twists; their additional
finite cyclotomic factors have \(2\)-power order.

On a geometric fibre, the finite \(H^0\) groups form a
Mittag--Leffler system. The continuous-cohomology comparison
\cite[Sections~18--21]{Milne} consequently gives
\[
 H^1(U,\mathcal G_r)
   =\left(\varprojlim_m H^1(U,\mathcal G_{r,m})\right)
       \otimes_{\mathcal O}E.
\]
The finite-level specialization isomorphisms therefore pass to the
\(E\)-adic groups. This argument does not require finite-level
cohomology to be free.

For compact supports we use relative duality, rather than applying
smooth proper base change to the open curve. The morphism \(g:U\to S\)
is smooth, separated, compactifiable, and of relative dimension one;
\(S\) is quasi-compact and quasi-separated, and every \(2^m\) is
invertible on it. The relative trace and smooth-duality theorems
\cite[Expose~XVIII, Theorem~2.9, Var.~2, and Theorem~3.2.5]{SGA4}
apply. Their base-change compatibility, followed by passage to \(E\),
identifies
\[
 H_c^1(U,\mathcal G_r)
       \simeq H^1(U,\mathcal G_r^\vee(1))^\vee
\]
compatibly with specialization. Moreover the natural transformation
\(Rg_!\to Rg_*\) is compatible with these base-change maps.
Accordingly the specialization square with horizontal maps
\(H_c^1\to H^1\) commutes, and both vertical maps are isomorphisms.
Taking images over \(E\) gives an isomorphism of the parabolic
cohomology groups, equivariant for residue-field Galois action.

The compatible finite systems give a continuous representation factoring
through \(\pi_1(S)\). Their parabolic image is a stable subquotient,
so \(H_r\) is unramified outside \(2,3\). At a good prime its
Frobenius is identified with the Frobenius of the special-fibre
parabolic group. This proves precisely the specialization needed for
the arithmetic trace calculation.

\subsection{The two trace formulas}

Let \(\alpha,\beta\) be the geometric Frobenius eigenvalues of a
smooth Legendre fibre over \(\F_q\). Thus
\(\alpha+\beta=a_q(\lambda)\) and \(\alpha\beta=q\). The
symmetric-power recurrence gives
\[
 \Tr(\Frob_\lambda\mid\Sym^2\mathcal V)
      =a_q(\lambda)^2-q,
 \qquad
 \Tr(\Frob_\lambda\mid\Sym^4\mathcal V)
      =a_q(\lambda)^4-3qa_q(\lambda)^2+q^2.
\]
The compact-support trace formula \cite{StacksTrace}, the vanishings
\eqref{eq:common-vanishing}, and the arithmetic boundary sequence
\eqref{eq:common-left-boundary} therefore give
\[
 \Tr(\Frob_q\mid H_r)
   =-\sum_{\lambda\ne0,1}
          \Tr(\Frob_\lambda\mid\mathcal G_{r,\lambda})-1.
\]
Since \(\sum_{\lambda\ne0,1}\rho_q(\lambda)=-1\), this is
\eqref{eq:trace-W2} for \(r=2\), and \eqref{eq:trace-V4} for
\(r=4\). The final \(-1\) is the trace of the actual boundary
line, not a correction inferred from its dimension. The preceding
specialization identifies these special-fibre traces with the
\(G_K\)-representations in Theorem~\ref{thm:common-geometry}, completing
its proof.

For powers of a fixed Frobenius, the character on \(\F_{q^d}^\times\)
must be \(\rho_q\circ N_{\F_{q^d}/\F_q}\). With this convention the
same formulas compute \(\Tr(\Frob_{\mathfrak p}^{d}\mid H_r)\).
At an inert rational prime \(p\equiv2\pmod3\), the initial residue
field is \(\F_{p^2}\); its cubic character is not a nontrivial norm
lift from \(\F_p\). These formulas determine the trace over \(K\)
and do not yet choose a trace at \(p\) for an extension to \(G_\Q\).

Parameter inversion gives
\(a_q(1/\lambda)=\chi_q(\lambda)a_q(\lambda)\), because the two
elliptic curves differ by the indicated quadratic twist. For even
\(k\), changing variables in the moment consequently yields
\[
 M_{3,k}(q;\rho_q)=M_{3,k}(q;\rho_q^{-1})\in\Z.
\]
Indeed the sum is a cyclotomic integer fixed by conjugation. This
character independence does not itself supply coefficient descent.

Section~\ref{sec:descent} places both cubic eigenspaces in an ambient
\(\Q_2\)-representation \(H_{\mathrm{nt}}\), constructs the
arithmetic inversion operator \(I\), and uses it to define the two
extensions \(W_4^\pm\) to \(G_\Q\).
